% =====================================================================
%  Phan Phuong phap, trinh bay theo mau Springer LNCS.
%
%  KHAC MOT DIEM SO VOI samplepaper.tex: mau goc dung [T1]{fontenc},
%  ma T1 KHONG co bang ma tieng Viet -- moi dau thanh se hong. Neu bai
%  viet bang tieng Viet thi phai dung [T5]{fontenc} nhu duoi day, hoac
%  chuyen sang XeLaTeX voi fontspec. Neu nop cho mot hoi nghi Springer
%  that thi bai phai bang tieng Anh, va luc do tra [T1] lai nhu cu.
%
%  Muc 1 va Muc 2 de trong co y, chi de so mucchay dung: phan Phuong
%  phap ra dung Muc 3, giong bai goc.
% =====================================================================
\documentclass[runningheads]{llncs}

\usepackage[utf8]{inputenc}
\usepackage[T5]{fontenc}
\usepackage{graphicx}
\usepackage{amsmath}
\usepackage{amssymb}
\usepackage{tikz}
\usetikzlibrary{arrows.meta, positioning}

\begin{document}

\title{Hợp nhất hai tầng cho định tuyến và phân lớp\\
trong học liên tục dựa trên tinh chỉnh hiệu quả tham số}

\author{Nguyễn Thiên Trường \and Tống Quang Thái}
\authorrunning{N. T. Trường và T. Q. Thái}
\institute{Trường Đại học Công nghệ, Đại học Quốc gia Hà Nội}

\maketitle

\begin{abstract}
Phần tóm tắt viết sau, khi đã chốt số liệu.

\keywords{Học liên tục \and Tinh chỉnh hiệu quả tham số \and Định tuyến
tham số \and Chiếu ngẫu nhiên.}
\end{abstract}

\section{Giới thiệu}

\emph{Để trống. Nội dung lấy từ báo cáo kỹ thuật.}

\section{Công trình liên quan}

\emph{Để trống.}

\section{Phương pháp}

\subsection{Phát biểu bài toán}

Học liên tục theo lớp (class-incremental learning, CIL) học một chuỗi $T$ tác vụ
mà không quên tác vụ cũ. Tác vụ thứ $t$ gồm $n_t$ mẫu
$\mathcal{D}_t=\{(x^t_i,y^t_i)\}_{i=1}^{n_t}$ thuộc tập lớp $\mathcal{C}_t$,
trong đó $x^t_i \in \mathcal{X}$ là ảnh vào và $y^t_i \in \mathcal{Y}$ là nhãn
tương ứng. Toàn bộ chuỗi là $\mathcal{D}=\{\mathcal{D}_1,\dots,\mathcal{D}_T\}$.

Ba ràng buộc định nghĩa thiết lập. Thứ nhất, tại thời điểm $t$ chỉ
$\mathcal{D}_t$ khả dụng cho huấn luyện. Thứ hai, các tập lớp không chồng lấn,
$\mathcal{C}_t \cap \mathcal{C}_{t'} = \emptyset$ với mọi $t \neq t'$, và khi suy
luận thì \emph{định danh tác vụ không được biết}: mô hình phải phân biệt mọi lớp
đã gặp bằng một đầu phân lớp dùng chung. Thứ ba, báo cáo này làm việc trong thiết
lập không lưu mẫu (rehearsal-free): ảnh của các tác vụ đã qua không bao giờ được
giữ lại, kể cả một tập nhỏ.

Trong các phương pháp dựa trên tinh chỉnh hiệu quả tham số (PET), mạng xương sống
tiền huấn luyện $h_{\mathrm{ptm}}$ với tham số $\theta_{\mathrm{ptm}}$ được đóng
băng, và một kho tham số nhẹ $\mathcal{P}=\{p_1,\dots,p_T\}$ được duy trì, mỗi bộ
huấn luyện riêng cho một tác vụ. Ở đây $p_t$ là một bộ điều hợp LoRA. Đầu phân
lớp dùng chung ký hiệu $g_{\theta_g}$. Vì định danh tác vụ không biết khi suy
luận, hệ thống buộc phải \emph{chọn} bộ tham số nào gắn vào mạng xương sống trước
khi phân lớp --- và toàn bộ chất lượng phía sau phụ thuộc vào lựa chọn đó.

\begin{table}
\caption{Ký hiệu dùng trong mục này.}\label{tab:notation}
\begin{tabular}{ll}
\hline
Ký hiệu & Ý nghĩa\\
\hline
$\mathcal{D}_t$, $\mathcal{C}_t$ & dữ liệu và tập lớp của tác vụ $t$\\
$h_{\mathrm{ptm}}$, $\theta_{\mathrm{ptm}}$ & mạng xương sống tiền huấn luyện, đóng băng\\
$\mathcal{P}=\{p_1,\dots,p_T\}$ & kho tham số LoRA, mỗi tác vụ một bộ\\
$g_{\theta_g}$, $g_\omega$ & đầu phân lớp dùng chung; bộ phân lớp phụ của TII\\
$q(x)$ & đặc trưng truy vấn từ mạng xương sống trần\\
$\hat{d}(x)$, $\tilde{d}(x)$ & phân bố theo lớp của TII; của đầu phân lớp sau định tuyến\\
$\mathcal{T}(\cdot)$ & ánh xạ lớp sang định danh tác vụ\\
$\hat{t}_f$, $\hat{t}_s$ & định danh tác vụ ở đối sánh đầu và sau khi đối sánh lại\\
$\ell \in \mathbb{R}^{|\mathcal{Y}|}$ & điểm số theo lớp của đầu phân lớp đã định tuyến\\
$s \in \mathbb{R}^{|\mathcal{Y}|}$ & điểm số theo lớp của đầu chiếu ngẫu nhiên\\
$z(\cdot)$ & chuẩn hóa về trung bình không, phương sai đơn vị trên từng mẫu\\
\hline
\end{tabular}
\end{table}

\subsection{Quy trình cơ sở}

Chúng tôi lấy HRM-PET làm cơ sở. Quy trình suy luận của nó gồm bốn khối nối
tiếp, trình bày dưới đây theo đúng thứ tự dữ liệu đi qua.

\subsubsection{Suy luận định danh tác vụ (TII).}
Ảnh vào được đưa qua mạng xương sống \emph{trần}, không gắn bộ điều hợp nào, để
lấy đặc trưng truy vấn. Đặc trưng này đi vào một bộ phân lớp phụ $g_\omega$ được
huấn luyện bằng entropy chéo, sinh ra một phân bố theo lớp:
\begin{equation}
q(x) = h(x;\theta_{\mathrm{ptm}}),
\qquad
\hat{d}(x) = g\!\left(q(x); \theta_\omega\right).
\label{eq:tii}
\end{equation}
Định danh tác vụ ban đầu thu được bằng cách lấy lớp có xác suất cao nhất rồi ánh
xạ sang tác vụ chứa nó, với $\mathcal{T}$ là ánh xạ lớp sang tác vụ.

\subsubsection{Đối sánh lại trực tiếp (DRM).}
Nhờ tri thức dùng chung trong kho tham số, với một số mẫu thì lớp mà đầu phân
lớp dự đoán \emph{đã} thuộc đúng tác vụ, dù đối sánh ban đầu sai. Những mẫu đó
nhận ra được vì lớp dự đoán không thuộc tác vụ đã đối sánh. HRM-PET so sánh hai
định danh:
\begin{align}
\hat{t}_f &= \mathcal{T}\!\left(\arg\max_i\, g(q(x);\theta_\omega)\right),
\label{eq:tf}\\
\hat{t}_s &= \mathcal{T}\!\left(\arg\max_i\,
   g\!\left(h(x; p_{\hat{t}_f}, \theta_{\mathrm{ptm}}); \theta_g\right)\right),
\label{eq:ts}
\end{align}
và chỉ thay thế khi hai giá trị khác nhau:
\begin{equation}
\hat{y} =
\begin{cases}
\arg\max_i\, g\!\left(h(x; p_{\hat{t}_s}, \theta_{\mathrm{ptm}}); \theta_g\right),
  & \text{nếu } \hat{t}_f \neq \hat{t}_s,\\[2pt]
\arg\max_i\, g\!\left(h(x; p_{\hat{t}_f}, \theta_{\mathrm{ptm}}); \theta_g\right),
  & \text{ngược lại.}
\end{cases}
\label{eq:drm}
\end{equation}
Nếu $\hat{t}_s$ đúng thì độ chính xác định tuyến tăng; nếu sai thì kết quả xấu
nhất cũng chỉ là dự đoán sai như $\hat{t}_f$ vốn đã sai. Thiết kế do đó không thể
làm giảm hiệu năng.

\subsubsection{Đối sánh lại theo độ tin cậy (CRM).}
Khi $\hat{t}_f = \hat{t}_s$ thì DRM không có gì để thay. Với những mẫu đó, HRM-PET
dựa trên một quan sát khác: mỗi bộ tham số chỉ được huấn luyện cùng các lớp của
tác vụ nó, nên với mẫu bị ghép sai, phân bố dự đoán thường kém chắc chắn hơn. Độ
tin cậy $E(\cdot)$ được tính bằng hàm entropy suy rộng hậu nghiệm, không cần huấn
luyện thêm, và một bộ phát hiện so nó với ngưỡng $\tau$:
\begin{equation}
De(x) =
\begin{cases}
0, & \text{nếu } E(\tilde{d}(x)) \leq \tau \quad (\text{nghi ghép sai}),\\
1, & \text{ngược lại.}
\end{cases}
\label{eq:detect}
\end{equation}
Vì không thể vạch một ranh giới dứt khoát bằng $\tau$, HRM-PET không thay thẳng
$\hat{t}_f$ mà so độ tin cậy giữa một tập ứng viên $\Gamma$ gồm $N$ lớp có xác
suất cao nhất trong $\hat{d}(x)$, rồi chọn định danh cho độ tin cậy lớn nhất:
\begin{equation}
\hat{t}_s = \arg\max_{i \,\in\, \Gamma}\;
   E\!\left(g\!\left(h(x; p_{\mathcal{T}(i)}, \theta_{\mathrm{ptm}});
   \theta_g\right)\right),
\qquad N = 2 .
\label{eq:crm}
\end{equation}
Mỗi ứng viên tốn thêm một lượt truyền xuôi, nên $N$ được giữ nhỏ.

\subsubsection{Đầu phân lớp dùng chung.}
Bộ tham số đã chọn được gắn vào mạng xương sống, và đầu phân lớp dùng chung sinh
véc-tơ điểm số theo lớp trên toàn bộ các lớp đã gặp:
\begin{equation}
\ell = g\!\left(h(x; p_{\hat{t}}, \theta_{\mathrm{ptm}}); \theta_g\right)
   \in \mathbb{R}^{|\mathcal{Y}|},
\qquad
\hat{y} = \arg\max_c\, \ell_c ,
\label{eq:logits}
\end{equation}
trong đó $\hat{t}$ là định danh tác vụ sau DRM và CRM. Các lớp chưa gặp bị che
bằng $-\infty$ nên không thể được dự đoán.

\subsubsection{Điểm yếu của cấu trúc này.}
Hai nhận xét dẫn tới phương pháp đề xuất. Thứ nhất, cả chuỗi bốn khối trên phụ
thuộc vào \emph{một} nguồn tín hiệu định tuyến duy nhất, sinh ra ở
(\ref{eq:tii}); DRM và CRM chỉ tinh chỉnh chứ không bổ sung nguồn nào. Thứ hai,
$\ell$ ở (\ref{eq:logits}) là đầu phân lớp \emph{dùng chung}, nên nó thường đã
đoán đúng lớp bất chấp định tuyến sai --- điều đó có nghĩa cải thiện định tuyến
không tự động chuyển thành cải thiện phân lớp.

Mục này đề xuất một nguồn tín hiệu thứ hai và hai vị trí để hợp nhất nó vào, ứng
với hai nhận xét trên. Cả hai đặt ở giai đoạn suy luận nên không huấn luyện lại
bất kỳ tham số nào. Hình~\ref{fig:pipeline} minh họa vị trí của chúng.

\subsection{Nguồn tín hiệu thứ hai}

Nguồn thứ hai phải thỏa hai ràng buộc. Thứ nhất, nó không được thêm tham số
huấn luyện, vì bất kỳ tham số nào cập nhật bằng gradient qua chuỗi tác vụ đều có
thể bị tác vụ sau ghi đè. Thứ hai, nó phải giữ được giao thức không lưu mẫu.

Chúng tôi dùng một đầu phân lớp bậc hai theo kiểu RanPAC. Đặc trưng $f$ được
trích từ \emph{một} bộ tham số LoRA cố định của HRM-PET, đóng vai trò thích nghi
phiên đầu. Đặc trưng được chiếu lên không gian số chiều cao bằng một ma trận
ngẫu nhiên đóng băng $W$, đi qua một hàm phi tuyến, rồi được tổng hợp thành
thống kê bậc hai:
\begin{align}
h &= \mathrm{ReLU}(f W), \qquad W \text{ đóng băng, sinh lại từ seed},
\label{eq:proj}\\
G &= \sum_n h_n h_n^{\top}, \qquad
C = \sum_n h_n\,\mathrm{onehot}(y_n)^{\top}.
\label{eq:gram}
\end{align}
Điểm số theo lớp là nghiệm đóng của hồi quy ridge trên hai thống kê đó:
\begin{equation}
s = h^{\top} (G + \lambda I)^{-1} C .
\label{eq:ridge}
\end{equation}

Ba nhận xét về thiết kế này. Ma trận chiếu ngẫu nhiên và đóng băng không phải
một sự đánh đổi mà là điều kiện cần: sinh lại được từ một seed cố định nên không
cần lưu trữ, và không bao giờ thay đổi nên không thể bị ghi đè. Hàm phi tuyến là
bắt buộc, vì thiếu nó thì phép chiếu chỉ là một biến đổi tuyến tính và không
thêm được khả năng tách lớp nào. Nghịch đảo ma trận Gram khử tương quan giữa các
chiều đã mở rộng, việc mà hiệp phương sai theo từng lớp trong không gian gốc
không làm được.

Quan trọng nhất, $G$ và $C$ ở (\ref{eq:gram}) đều là \emph{tổng} trên các mẫu.
Phép cộng có tính giao hoán, nên kết quả không phụ thuộc thứ tự các tác vụ được
học: đầu phân lớp này không quên theo đúng nghĩa cấu trúc, chứ không phải nhờ một
cơ chế chống quên nào. Và vì chỉ giữ hai ma trận tổng hợp, không giữ đặc trưng
của từng mẫu, giao thức không lưu mẫu được bảo toàn.

\subsection{Hợp nhất hai tầng}

\subsubsection{Tầng thứ nhất: hợp nhất ở cấp định tuyến.}
Cả TII lẫn đầu chiếu ngẫu nhiên đều sinh điểm số theo lớp, trong khi định tuyến
cần điểm số theo tác vụ. Chúng tôi quy đổi bằng cách lấy giá trị lớn nhất trên
tập lớp của mỗi tác vụ:
\begin{equation}
\ell^{\mathrm{task}}_t = \max_{c \,\in\, \mathcal{C}_t} \ell_c ,
\qquad
s^{\mathrm{task}}_t = \max_{c \,\in\, \mathcal{C}_t} s_c .
\label{eq:taskscore}
\end{equation}
Hai véc-tơ này nằm trên hai thang đo khác nhau: điểm ridge trải khoảng một đơn
vị trong khi logit trải khoảng ba mươi lăm. Trộn thẳng sẽ để TII lấn át ở mọi
trọng số, nên mỗi nguồn được chuẩn hóa về trung bình không và phương sai đơn vị
trên từng mẫu bằng toán tử $z(\cdot)$ trước khi trộn:
\begin{equation}
\mathrm{fused} = w\, z(\ell^{\mathrm{TII}}) + (1-w)\, z(s^{\mathrm{RP}}),
\qquad w = 0.7 .
\label{eq:route}
\end{equation}
Định danh tác vụ $\arg\max_t \mathrm{fused}_t$ được trao cho DRM và CRM như một
điểm khởi đầu. Chúng tôi thay \emph{đúng một} thứ là bước đối sánh đầu tiên; mọi
cơ chế phía sau của HRM-PET giữ nguyên. Tại $w = 1.0$, biểu
thức~(\ref{eq:route}) trả về đúng thứ tự của $\ell^{\mathrm{TII}}$, nên phương
pháp thu về baseline một cách chính xác.

\subsubsection{Tầng thứ hai: hợp nhất ở cấp phân lớp với cơ chế cổng.}
Định tuyến tốt hơn không tự động thành phân lớp tốt hơn. Đầu phân lớp dùng chung
thường đã đoán đúng lớp bất chấp định tuyến sai, nên sửa định tuyến ở những mẫu
đó không đổi được dự đoán. Chúng tôi vì vậy dùng chính đầu chiếu ngẫu nhiên làm
bộ phân lớp thứ hai.

Một trọng số cố định là sai. Nó chi cùng một phần quyết định cho mẫu mà đầu
chính gọi tên ở mức $p = 0.97$ như cho mẫu nó đang phân vân, mà loại thứ nhất
chiếm đa số áp đảo; trả giá ở đó chỉ làm phẳng một đỉnh vốn đã đúng. Hợp nhất
chỉ đổi được dự đoán bằng cách lật đổ lớp đứng đầu, và đối thủ khả dĩ nhất của
nó là lớp đứng thứ hai. Do đó trọng số được đặt tỉ lệ nghịch với biên giữa hai
lớp cao nhất của chính đầu chính:
\begin{equation}
\beta_i = \beta \left( 1 - \frac{p_{(1)} - p_{(2)}}{p_{(1)} + p_{(2)}} \right),
\qquad \beta = 0.5 ,
\label{eq:gate}
\end{equation}
trong đó $p_{(1)}$ và $p_{(2)}$ là hai xác suất lớn nhất của đầu chính. Chia cho
tổng chứ không phải cho một khiến đại lượng này chỉ phụ thuộc tỉ số
$p_{(2)}/p_{(1)}$, tức đo mức độ lớp đứng đầu bị đe dọa, độc lập với khối lượng
xác suất mà cặp đó nắm giữ.

Phép trộn dùng lại phép chuẩn hóa của tầng thứ nhất, rồi đưa kết quả về đúng
thang đo của $\ell$:
\begin{align}
m_i &= (1 - \beta_i)\, z(\ell_i) + \beta_i\, z(s_i),
\label{eq:mix}\\
\hat{\ell}_i &= m_i \cdot \mathrm{std}(\ell_i) + \mathrm{mean}(\ell_i).
\label{eq:affine}
\end{align}
Bước đưa về ở~(\ref{eq:affine}) là một phép biến đổi affine theo từng mẫu, nên
thứ hạng và do đó độ chính xác không đổi; nó chỉ khiến hàm mất mát so sánh được
với baseline, vì chuẩn hóa đặt hỗn hợp ở phương sai đơn vị, một thang đo sai cho
entropy chéo. Tại $\beta = 0$ thì $\hat{\ell} = \ell$ đúng từng phần tử, nên
tầng này cũng thu về baseline một cách chính xác.

\begin{figure}
\centering
\resizebox{\textwidth}{!}{%
\begin{tikzpicture}[
  font=\footnotesize,
  >={Stealth[length=2mm]},
  box/.style={draw, minimum height=8mm, minimum width=1.6cm, inner xsep=4pt, align=center},
  hl/.style={box, very thick, fill=black!6}
]
\node[box] (in) at (0.0, 0.0)    {Ảnh vào};
\node[box] (tii) at (2.7, 0.85)  {TII};
\node[box] (rp) at (2.7, -0.85)  {Đầu RP};
\node[hl]  (f1) at (5.9, 0.0)    {Tầng 1\\Hợp nhất định tuyến};
\node[box] (lo) at (8.9, 0.0)    {Chọn LoRA};
\node[box] (dr) at (11.3, 0.0)   {DRM / CRM};
\node[box] (cl) at (13.9, 0.0)   {Đầu phân lớp};
\node[hl]  (f2) at (13.9, -2.1)  {Tầng 2\\Hợp nhất phân lớp};
\node[box] (out) at (13.9, -3.7) {Dự đoán};
\draw[->] (in) -- (tii);
\draw[->] (in) -- (rp);
\draw[->] (tii) -- (f1);
\draw[->] (rp) -- (f1);
\draw[->] (f1) -- (lo);
\draw[->] (lo) -- (dr);
\draw[->] (dr) -- (cl);
\draw[->] (cl) -- (f2);
\draw[->] (f2) -- (out);
\draw[->, dashed] (rp.south) |- node[pos=0.75, above, font=\scriptsize]
      {điểm số theo lớp của đầu RP} (f2.west);
\end{tikzpicture}}
\caption{Quy trình suy luận. Các khối viền đậm là hai thành phần được bổ sung.
Đầu chiếu ngẫu nhiên được tính một lần duy nhất và cấp tín hiệu cho cả hai tầng:
điểm số theo tác vụ cho tầng thứ nhất, điểm số theo lớp cho tầng thứ hai (đường
nét đứt).}
\label{fig:pipeline}
\end{figure}

\subsubsection{Chi phí.}
Đầu chiếu ngẫu nhiên cần đúng một lượt truyền xuôi trên mỗi mẫu, và cùng một
kết quả đó phục vụ cả hai tầng. Không tham số nào được huấn luyện thêm, và không
checkpoint nào phải huấn luyện lại.

\end{document}
